% ========== 3.1 SYSTEM REQUIREMENTS ANALYSIS ==========

\clearpage
\section{3.1 System Requirements Analysis}
\addcontentsline{toc}{section}{3.1 System Requirements Analysis}
\label{sec:system-requirements-analysis}

\lettrine{T}{his section} analyzes how the requirements identified in Chapter 2 translate into architectural decisions and design constraints. Rather than repeating the requirements themselves, we examine their implications for system architecture, identify trade-offs between competing needs, and establish the technical foundation for design choices that follow.

\subsection*{3.1.1 Architectural Implications of Requirements}

The functional requirements for AI orchestration demand a sophisticated coordination engine capable of managing multiple concurrent agents. This drives the design of the Conductor component with reinforcement learning capabilities. The extension system requirements necessitate robust plugin architecture with secure sandboxing, leading to the Dual Ensemble Architecture that separates trusted infrastructure from community extensions.

Performance requirements establish stringent targets: rapid startup time eliminates monolithic architectures in favor of minimal microkernel design. Minimal memory footprint necessitates lazy loading and efficient resource management. These constraints drive the decision to implement infrastructure in Rust for optimal efficiency.

Security requirements demand process isolation for untrusted extensions, capability-based access control, and comprehensive input validation. The need to support community extensions while maintaining security leads to the sandboxed execution model of The Grand Stage with monitored resource usage.

\subsection*{3.1.2 Design Trade-offs and Priorities}

Key trade-offs shape architectural decisions. Performance versus flexibility: in-process execution (The Pit) provides microsecond latency for trusted components, while out-of-process execution (The Grand Stage) adds millisecond overhead but ensures safety. Security versus convenience: strict sandboxing protects users but limits extension capabilities, resolved through graduated trust levels and Player registration.

Scalability versus simplicity: supporting projects from small scripts to enterprise codebases requires sophisticated resource management, balanced through modular architecture where users enable features based on needs. The actor-based concurrency model and stateless component design enable horizontal scaling for cloud deployments.

\subsection*{3.1.3 Quality Validation Criteria}

Quality attribute scenarios provide measurable validation criteria. Performance: "Large projects become responsive quickly" and "Code completion appears with minimal latency in most cases." Security: "Malicious extension access attempts are blocked and logged" and "Extension crashes don't affect other components."

Usability: "New users create AI-assisted code within minutes" and "Error messages provide actionable guidance in most cases." Maintainability: "New AI model integration implemented efficiently" and "Bug location identified quickly using documentation."

These scenarios drive testing efforts and validate that architectural decisions satisfy stakeholder needs while maintaining essential quality attributes.


